% ══════════════════════════════════════════════════════════════════
%  FICHE D'EXERCICES — Chapitre 9 : Les acides α-aminés et les protéines
% ══════════════════════════════════════════════════════════════════
\section{Exercices type Baccalauréat}
\textit{Données : $M(\ce{H})=1$, $M(\ce{C})=12$, $M(\ce{N})=14$, $M(\ce{O})=16$
\si{\gram\per\mol}.}

\rubrique{Structure, nomenclature}

\exo{}\diff{1} Donner la formule générale d'un acide $\alpha$-aminé. Identifier les
deux groupes fonctionnels caractéristiques.

\exo{}\diff{1} La glycine a pour formule $\ce{H2N-CH2-COOH}$ et l'alanine
$\ce{CH3-CH(NH2)-COOH}$. Calculer leur masse molaire et entourer, pour l'alanine, le
carbone asymétrique.

\rubrique{Chiralité}

\exo{}\diff{2} Montrer que l'alanine possède un atome de carbone asymétrique.
Représenter ses deux énantiomères et expliquer pourquoi la glycine, elle, n'est pas chirale.

\exo{}\diff{2} Combien de stéréo-isomères présente un acide $\alpha$-aminé possédant
deux carbones asymétriques ? Justifier.

\rubrique{Caractère amphotère}

\exo{}\diff{2} En solution aqueuse, un acide $\alpha$-aminé existe principalement sous
forme d'ion dipolaire (zwitterion).
\begin{enumerate}[label=\arabic*)]
\item Écrire la formule du zwitterion de la glycine.
\item Écrire les équations traduisant le comportement de la glycine en milieu acide
puis en milieu basique.
\item Définir le pH isoélectrique (pH$_i$).
\end{enumerate}

\exo{}\diff{2} Pour la glycine, on donne $pK_{a1} = 2{,}3$ (couple
$\ce{H3N+-CH2-COOH}/\ce{H3N+-CH2-COO-}$) et $pK_{a2} = 9{,}6$ (couple
$\ce{H3N+-CH2-COO-}/\ce{H2N-CH2-COO-}$). Déterminer la forme prédominante de la
glycine à pH $= 1$, à pH $= 6$ et à pH $= 12$.

\rubrique{Liaison peptidique}

\exo{}\diff{1} Définir la liaison peptidique. Écrire l'équation de formation du
dipeptide obtenu à partir de deux molécules de glycine, et entourer la liaison peptidique.

\exo{}\diff{2} On condense une molécule de glycine et une molécule d'alanine.
\begin{enumerate}[label=\arabic*)]
\item Montrer qu'il existe deux dipeptides possibles. Les représenter.
\item Quelle est la différence entre ces deux dipeptides ?
\end{enumerate}

\rubrique{Détermination}

\exo{}\diff{3} Un acide $\alpha$-aminé $A$ a pour formule brute $\ce{C3H7O2N}$.
\begin{enumerate}[label=\arabic*)]
\item Calculer sa masse molaire.
\item Écrire sa formule semi-développée (acide $\alpha$-aminé) et le nommer.
\item $A$ est-il chiral ? Justifier.
\item Écrire la formule de son zwitterion et les équations de ses réactions avec
$\ce{H3O+}$ et avec $\ce{HO-}$.
\end{enumerate}

\rubrique{Problème de synthèse — type Baccalauréat}

\sujetbac[d'après Baccalauréat D, Cameroun]
On étudie un acide $\alpha$-aminé naturel $A$ qui contient en masse \SI{15.7}{\percent}
d'azote. Sa molécule ne possède qu'une fonction acide et qu'une fonction amine.
\begin{enumerate}[label=\arabic*)]
\item En supposant que $A$ contient un seul atome d'azote, déterminer sa masse molaire.
\item Sachant que $A$ est saturé et de formule $\ce{C_xH_yO2N}$, déterminer $x$ et $y$.
\item Écrire la (les) formule(s) semi-développée(s) possible(s) d'acide
$\alpha$-aminé et la (les) nommer.
\item Indiquer celle(s) qui présente(nt) une activité optique (chiralité).
\end{enumerate}
